%\documentclass[landscape,a0b,final,a4resizeable]{a0poster}
\documentclass[landscape,a0b,final]{a0poster}
%\documentclass[portrait,a0b,final,a4resizeable]{a0poster}
%\documentclass[portrait,a0b,final]{a0poster}
%%% Option "a4resizeable" makes it possible ot resize the
%   poster by the command: psresize -pa4 poster.ps poster-a4.ps
%   For final printing, please remove option "a4resizeable" !!

\usepackage{epsfig}
\usepackage{multicol}
\usepackage{pstricks,pst-grad}

%%%%%%%%%%%%%%%%%%%%%%%%%%%%%%%%%%%%%%%%%%%
% Definition of some variables and colors
%\renewcommand{\rho}{\varrho}
%\renewcommand{\phi}{\varphi}
\setlength{\columnsep}{3cm}
\setlength{\columnseprule}{2mm}
\setlength{\parindent}{0.0cm}

%My commands
\usepackage{overcite}
\usepackage{amsmath}
\usepackage{amsfonts}
\usepackage{bm}
\newcommand{\mat}[1]{\mathbb{#1}} % mtpro special fonts
\newcommand{\tensortwo}[1]{\mat{#1}}
\newcommand{\bindent}{\ \ \ \ }

%%%%%%%%%%%%%%%%%%%%%%%%%%%%%%%%%%%%%%%%%%%%%%%%%%%%
%%%               Background                     %%%
%%%%%%%%%%%%%%%%%%%%%%%%%%%%%%%%%%%%%%%%%%%%%%%%%%%%

\newcommand{\background}[3]{
  \newrgbcolor{cgradbegin}{#1}
  \newrgbcolor{cgradend}{#2}
  \psframe[fillstyle=gradient,gradend=cgradend,
  gradbegin=cgradbegin,gradmidpoint=#3](0.,0.)(1.\textwidth,-1.\textheight)
}

%%%%%%%%%%%%%%%%%%%%%%%%%%%%%%%%%%%%%%%%%%%%%%%%%%%%
%%%                Poster                        %%%
%%%%%%%%%%%%%%%%%%%%%%%%%%%%%%%%%%%%%%%%%%%%%%%%%%%%

\newenvironment{poster}{
  \begin{center}
  \begin{minipage}[c]{0.98\textwidth}
}{
  \end{minipage} 
  \end{center}
}

%%%%%%%%%%%%%%%%%%%%%%%%%%%%%%%%%%%%%%%%%%%%%%%%%%%%
%%%                pcolumn                       %%%
%%%%%%%%%%%%%%%%%%%%%%%%%%%%%%%%%%%%%%%%%%%%%%%%%%%%

\newenvironment{pcolumn}[1]{
  \begin{minipage}{#1\textwidth}
  \begin{center}
}{
  \end{center}
  \end{minipage}
}

%%%%%%%%%%%%%%%%%%%%%%%%%%%%%%%%%%%%%%%%%%%%%%%%%%%%
%%%                pbox                          %%%
%%%%%%%%%%%%%%%%%%%%%%%%%%%%%%%%%%%%%%%%%%%%%%%%%%%%

\newrgbcolor{lcolor}{0. 0. 0.80}
\newrgbcolor{gcolor1}{1. 1. 1.}
\newrgbcolor{gcolor2}{.80 .80 1.}

\newcommand{\pbox}[4]{
\psshadowbox[#3]{
\begin{minipage}[t][#2][t]{#1}
#4
\end{minipage}
}}

%%%%%%%%%%%%%%%%%%%%%%%%%%%%%%%%%%%%%%%%%%%%%%%%%%%%
%%%                myfig                         %%%
%%%%%%%%%%%%%%%%%%%%%%%%%%%%%%%%%%%%%%%%%%%%%%%%%%%%
% \myfig - replacement for \figure
% necessary, since in multicol-environment 
% \figure won't work

\newcommand{\myfig}[3][0]{
\begin{center}
  \vspace{0.25cm}
  \includegraphics[width=#3\hsize,angle=#1]{#2}
  \nobreak\medskip
\end{center}}

%%%%%%%%%%%%%%%%%%%%%%%%%%%%%%%%%%%%%%%%%%%%%%%%%%%%
%%%                mycaption                     %%%
%%%%%%%%%%%%%%%%%%%%%%%%%%%%%%%%%%%%%%%%%%%%%%%%%%%%
% \mycaption - replacement for \caption
% necessary, since in multicol-environment \figure and
% therefore \caption won't work

%\newcounter{figure}
\setcounter{figure}{1}
\newcommand{\mycaption}[1]{
  \vspace{0.25cm}
  \begin{quote}
    {{\sc Figure} \arabic{figure}: #1}
  \end{quote}
  \vspace{0.25cm}
  \stepcounter{figure}
}

%%%%%%%%%%%%%%%%%%%%%%%%%%%%%%%%%%%%%%%%%%%%%%%%%%%%%%%%%%%%%%%%%%%%%%
%%% Begin of Document
%%%%%%%%%%%%%%%%%%%%%%%%%%%%%%%%%%%%%%%%%%%%%%%%%%%%%%%%%%%%%%%%%%%%%%

\begin{document}

\newrgbcolor{lightblue}{0. 0. 0.80}
\newrgbcolor{white}{1. 1. 1.}
\newrgbcolor{whiteblue}{.80 .80 1.}
\newrgbcolor{ubblue}{0.0157 0.102 0.608}
\newrgbcolor{ubgray}{0.4 0.4 0.4}

\background{1. 1. 1.}{1. 1. 1.}{0.5}

\vspace*{1cm} %Controls white space on top of poster

\begin{poster}

%%%%%%%%%%%%%%%%%%%%%
%%% Header
%%%%%%%%%%%%%%%%%%%%%
\begin{center}
\begin{pcolumn}{0.98}

\pbox{0.98\textwidth}{}{linewidth=2mm,framearc=0.3,linecolor=ubblue,fillstyle=gradient,gradangle=0,gradbegin=ubgray,gradend=white,gradmidpoint=1.0,framesep=1em}{
    %%% Unisiegel
    \begin{minipage}[c][9cm][c]{0.1\linewidth}
    \begin{center}
        \includegraphics[width=20cm,angle=0]{figures/ublogo}
    \end{center}
    \end{minipage}
    %%% Titel
    \begin{minipage}[c][9cm][c]{0.75\textwidth}
    \begin{center}
        %hspace commands control centering of title above poster. You may need to mess with these depending on how long your title is.
        \hspace{5cm}
        {\sc \huge Analysis of Optical Activity in Terms of Bond and Lone-Pairs:}\\[5mm] 
        \hspace{5cm}
        {\sc \huge The Exceptionally Large Optical Rotation of Norborneneone\cite{Autschbach:2012k}}\\[5mm]
        \hspace{5cm}
        {\Large Barry Moore II, Monika Srebro, and Jochen Autschbach\\[7.5mm]
        \hspace{5cm}
        University at Buffalo, State University of New York, Buffalo, NY, 14260}
    \end{center}
    \end{minipage}
    %%% GK-Logo
    \begin{minipage}[c][9cm][c]{0.01\linewidth}
    \begin{flushright}
        %This hspace puts the nsf symbol at the far right of the poster banner. May need to change this but don't remove the '*'
        \hspace*{5cm}
        \includegraphics[width=10cm,angle=0]{figures/nsf4}
    \end{flushright}
    \end{minipage}
}
\end{pcolumn}
\end{center}

\vspace*{0.1cm} %Spacing between banner and text

%%%%%%%%%%%%%%%%%%%%%
%%% Content
%%%%%%%%%%%%%%%%%%%%%

%%% Begin of Multicols-Enviroment
\begin{multicols}{3}

%%% Abstract
\begin{center}\pbox{0.8\columnwidth}{}{linewidth=2mm,framearc=0.3,linecolor=ubblue,fillstyle=gradient,gradangle=0,gradbegin=white,gradend=ubgray,gradmidpoint=1.0,framesep=1em}{\begin{center} \Large \color{ubblue}{\bf{Abstract}}\end{center}}\end{center} %Makes the Abstract banner 
\vspace{0.25cm}
\large \bindent Norbornenone, which has both a C$=$O and a C$=$C chromophore in a rigid bicyclic hydrocarbon framework,
exhibits optical rotation (OR) an order of magnitude larger than that of similar molecules with only one of these chromophores
(e.g., $\alpha-$pinene). Its OR is also very sensitive to approximations in electronic structure calculations. The present study
demonstrates a novel approach to interpret optical rotation using familiar concepts of chemical bonding, aided by first-principles
calculations. A theoretical procedure is developed for analyzing the OR tensor components of a molecule in terms of individual
bonds and lone pairs. The link between the chemist’s bond and quantum mechanics is provided by localized molecular orbitals
obtained from density functional theory (DFT) calculations. Delocalization of $\pi$ orbitals is shown to play a crucial role in the
large OR of norbornenone, as hinted by the DFT delocalization error inherent in many standard functionals and confirmed by
detailed analysis. The OR contributions generated by the double bond in $\alpha-$pinene are even stronger than that of norbornenone.
The isotropic average, observed in solution or in gas phase, is small as a result of cancellation of tensor components with opposite
signs. The electronic coupling and delocalization of the C$=$C $\pi$ bond and the C$=$O oxygen $\pi$ lone pair in norbornenone
selectively enhance one of the OR tensor components, resulting in the exceptionally large negative isotropic OR.

%%% Functional Tuning
\vspace{0.25cm}
\begin{center}\pbox{0.8\columnwidth}{}{linewidth=2mm,framearc=0.3,linecolor=ubblue,fillstyle=gradient,gradangle=0,gradbegin=white,gradend=ubgray,gradmidpoint=1.0,framesep=1em}{\begin{center} \Large \color{ubblue}{\bf{Functional Tuning}}\end{center}}\end{center}
\vspace{0.25cm}

    \bindent We use the following 3-parameter expression for the inverse
    interelectronic distance in the exchange functional
    \cite{yanai_new_2004}:
        \vspace{-1.8cm}
    \begin{equation}
        \vspace{-1.8cm}
    \label{eq:r12-range-separation-alpha-beta-gamma}
    \frac{1}{r_{12}} =
    \frac{\alpha + \beta \operatorname{erf} (\gamma r_{12})}{r_{12}} + 
    \frac{1-[\alpha + \beta \operatorname{erf} (\gamma r_{12})]}{r_{12}} 
    \end{equation}
    where the first term on the right-hand side is used for the long-range
    HF part. In Equation (\ref{eq:r12-range-separation-alpha-beta-gamma}),
    $\gamma$ is the range-separation parameter, $\alpha$ is the HF
    fraction at very small interelectronic distances, and $\alpha+\beta$
    quantifies the fraction of HF exchange in the asymptotic limit at
    large interelectronic separations. A global hybrid such as B3LYP
    corresponds to $\beta = 0$, with $\alpha$ being the HF fraction of
    exchange in the functional.

    \bindent We consider a range-separated hybrid variant of the PBE functional
    which satisfies the correct asymptotic behavior of the XC
    potential. It follows that $\alpha + \beta = 1$. Further, the
    functional is tuned to satisfy
    $-{\varepsilon}^{\text{HOMO}}(N)=\text{IP}(N)$ subject to the $\alpha
    + \beta = 1$ constraint. Following Reference
    \citen{kuritz_charge-transfer-like_2011}, we minimize
        \vspace{-1.8cm}
    \begin{equation}
        \vspace{-1.8cm}
    J^2 =
    \sum_{i=0}^{1} [\varepsilon^{\text{HOMO}}(N+i)+\text{IP}(N+i)]^2 
    \label{eq:J2-criterion}
    \end{equation} 
    to satisfy the IP condition simultaneously for the system with $N$ and
    $N+1$ electrons, where $\text{IP}(N)= E(N)-E(N-1)$ is calculated as a
    difference of total energies. This optimization is performed in the
    $\alpha = 1-\beta$ vs.\ $\gamma$ plane, which is the reason why the
    procedure is referred to as a `two-dimensional' tuning here. The
    tuning was accomplished as follows. Hartree-Fock molecular orbitals
    were calculated for the neutral ($N$ electrons), anion ($N+1$), and
    cation ($N-1$) systems. These molecular orbital vectors were then used
    as initial guesses for DFT calculations with the PBE0 global hybrid
    functional. The PBE0 molecular orbital vectors were used as input for
    calculations with $\alpha$ between 0 and 0.5 and $\gamma$ between 0
    and 0.5 on the tuning grid. MOs from calculations with $\alpha = 0.5$
    and a given $\gamma$ were used as initial guesses for calculations
    with $\alpha$ ranging from 0.6 to 0.8 and the same value of $\gamma$
    in order to improve convergence of the self-consistent field (SCF)
    procedure.  Finally, for sets $\{\alpha,\beta = 1-\alpha,\gamma\}$
    satisfying the IP condition, we investigate the DFT delocalization
    error, quantified by the curvature of $E(N)$.
    %%% Figures:
    \begin{center}
    \vspace*{1cm}
        % \myfig[rotation in degrees]{epsfile}{stretching parameter}
    \myfig[0]{figures/norbornenone-tuning-correct}{1.0}
    \mycaption{ a) Contour plot of $J^2$ as a function of
    $\alpha$ and $\gamma$. The thick orange line corresponds to $J^2
    \approx 0$. b) Energy of norbornenone as a function
    of fractional electron number, ${\Delta}N$, relative to neutral
    system $({\Delta}N = 0$) for tuned LC-PBE0-\# parameterizations. The
    numerical values in parentheses correspond to $({\Delta}N)^2$
    coefficients of quadratic fits to $E(N)$ in the electron-deficient
    and electron-rich regime, respectively $({\Delta}N < 0$, ${\Delta}N
    > 0)$. The inset shows the $E(N)$ behavior in the electron-rich
    regime.}
    \end{center}

%%% Analysis Protocol
\vspace{1.0cm}
\begin{center}
\pbox{0.8\columnwidth}{}{linewidth=2mm,framearc=0.1,linecolor=ubblue,fillstyle=gradient,gradangle=0,gradbegin=white,gradend=ubgray,gradmidpoint=1.0,framesep=1em}
{\begin{center} \Large \color{ubblue}{\bf{Analysis Protocol}}\end{center}}
\end{center}
\vspace{1.0cm}

    %%% Figures:
    \begin{center}
        % \myfig[rotation in degrees]{epsfile}{stretching parameter}
    \myfig[0]{figures/flowchart}{1.0}
    \mycaption{Flowchart for OR analysis. CNC = center of nuclear charges,
    PAS = principal axis system of $\tilde{\tensortwo{B}}$.}
    \end{center}
\vspace{5cm}    
%%% Results and Discussion
\vspace{0.25cm}
\begin{center}
\pbox{0.8\columnwidth}{}{linewidth=2mm,framearc=0.1,linecolor=ubblue,fillstyle=gradient,gradangle=0,gradbegin=white,gradend=ubgray,gradmidpoint=1.0,framesep=1em}
{\begin{center} \Large \color{ubblue}{\bf{Results and Discussion}}\end{center}}
\end{center}
\vspace{0.25cm}

    \bindent Electron delocalization plays a paramount role in the optical rotation
    (OR) of norbornenone. This conclusion is evident from an analysis in
    terms of individual `chemist's orbitals' (localized
    MOs). Delocalization of the C=C $\pi$ bond over the C=O group, and
    delocalization of the carbonyl oxygen $\pi$ lone-pair orbital over the
    C=C moiety, result in a significantly enhanced component of the OR
    tensor in the direction of the C=O bond, $\beta_{33}$. Contributions
    from the bridging C--C $\sigma$ bonds are also very important, as one
    might expect if delocalization between the C=C and C=O $\pi$ orbitals
    is significant. The strong enhancement of the negative $\beta_{33}$ is
    not balanced by other tensor components which results in a large
    negative isotropic optical rotation of norbornenone. A comparison with
    $\alpha$-pinene, which has a more localized C=C $\pi$ bond as a
    composition of the corresponding `natural LMO' revealed (magnitudes below $<$ 0.4
    on adjacent and distant carbons
    clearly shows how geometry and electronic structure effects in
    norbornenone all cooperate to yield an enhancement of $\beta_{33}$.

    \bindent It is clear that one can dig much deeper in an attempt to fully
    rationalize the magnitudes and the signs of all orbital contributions
    to the optical activity of norbornenone, $\alpha$-pinene and countless
    other molecules. The present work is intended to provide an initial
    assessment of the analysis technique and to address long standing
    problems with OR calculations for norbornenone. Further, we
    hope to have demonstrated that with suitable theoretical tools it is
    possible to gain an in-depth understanding of optical rotation. We
    expect to be able to provide additional examples in subsequent
    publications and, in particular, consider the optical rotation created
    by chiral $\sigma$-frameworks.
    %%% Figures:
    \begin{center}
        % \myfig[rotation in degrees]{epsfile}{stretching parameter}
    \myfig[0]{figures/norbornenone-pinene-comparison}{1.0}
    \mycaption{Static chiroptical response tensor $\tilde{\tensortwo{B}}$
        for norbornenone and $\alpha$-pinene, and principal
        components. Yellow and blue indicate negative and positive optical
        rotation for a given direction.  Surfaces were scaled to 30 pm per
        atomic unit. LC-PBE0* functional. Values in parentheses are
        diagonal components of $\tensortwo{\beta}$.}
    \end{center}

%%% Acknowledgements
\vspace{0.25cm}
\begin{center}
\pbox{0.8\columnwidth}{}{linewidth=2mm,framearc=0.1,linecolor=ubblue,fillstyle=gradient,gradangle=0,gradbegin=white,gradend=ubgray,gradmidpoint=1.0,framesep=1em}
{\begin{center} \Large \color{ubblue}{\bf{Acknowledgements}}\end{center}}
\end{center}
\vspace{0.25cm}
    
     This research has been supported by the National Science
     Foundation, Grant No. CHE 0952253 to J.A. M.S. is grateful
     for financial support from the Foundation for Polish Science
     (START scholarship) as well as from the Polish Ministry of
     Science and Higher Education (Mobility Plus program). We
     thank the Center for Computational Research at the University
     at Buffalo for supporting our research.
      
%%% References
\bibliographystyle{achemso}
\bibliography{autschbach,myjournals,misc,dft,cdrefs,response,cd}

\end{multicols}
\end{poster}
\end{document}

